\chapter{Independence, Basis, and Dimension}

By now you have met spans (Chapter~3), vector spaces and subspaces
(Chapter~11), and the two great subspaces attached to a matrix --- the column
space and the nullspace (Chapter~12). You have a feeling for what a subspace
\emph{is}: a flat sheet through the origin, a line, a plane, all of $\R^n$.

This chapter answers the next, deeper question. A subspace is built out of
vectors. \emph{How many vectors do you really need to build it, and how do you
know when you have exactly the right ones --- no wasted vectors, no missing
directions?} That is the whole story of this chapter, and it rests on three
words that depend on one another like rungs of a ladder:

\begin{itemize}
  \item \textbf{Independence} --- no vector in your set is redundant.
  \item \textbf{Spanning} --- your set is rich enough to reach everything.
  \item \textbf{Basis} --- the Goldilocks set: independent \emph{and} spanning, just right.
\end{itemize}

And out of ``basis'' falls a single number, the \textbf{dimension}, that pins
down the size of a space once and for all. By the end you will be able to look
at a column space or a nullspace and write down a clean, minimal set of vectors
that generates it, count them, and trust that count completely.

\chapterline

\section{The redundant vector: a picture of dependence}
\label{ch14:redundant}

Start with the simplest possible question. I hand you a pile of vectors. Are any
of them \emph{wasted}? A vector is wasted if you could have built it for free
out of the others --- if it adds no new direction to your reach.

\begin{intuitionbox}
Think of each vector as a road leaving a town at the origin. A set of roads is
\textbf{dependent} if at least one road takes you somewhere you could already
reach by combining the other roads. That road is redundant: tear it up and you
lose nothing. A set is \textbf{independent} if every single road opens up genuinely
new territory --- remove any one of them and some destination becomes
unreachable.
\end{intuitionbox}

Let us make that concrete in $\R^2$. Take
\[
  \vec{v}_1 = \colvec{2 \\ 1}, \qquad
  \vec{v}_2 = \colvec{-1 \\ 1}, \qquad
  \vec{v}_3 = \colvec{1 \\ 5}.
\]
The first two point in genuinely different directions: one leans right-and-up,
the other leans left-and-up. Neither is a scalar multiple of the other, so each
opens new territory. But watch the third vector:
\[
  2\,\vec{v}_1 + 3\,\vec{v}_2
  = 2\colvec{2\\1} + 3\colvec{-1\\1}
  = \colvec{4\\2} + \colvec{-3\\3}
  = \colvec{1\\5}
  = \vec{v}_3 .
\]
So $\vec{v}_3$ is not new at all. Anywhere $\vec{v}_3$ could take you, the pair
$\vec{v}_1, \vec{v}_2$ could already take you. The three vectors are
\textbf{dependent}: there is slack in the set, a vector you can throw away.

\begin{figure}[H]
\centering
\begin{tikzpicture}[>=Stealth, scale=0.95]
  % grid
  \draw[help lines, gray!30] (-4,-1) grid (5,6);
  \draw[->,gray!60] (-4,0)--(5,0) node[right]{$x$};
  \draw[->,gray!60] (0,-1)--(0,6) node[above]{$y$};
  % v1, v2
  \draw[->,very thick,vecblue] (0,0)--(2,1) node[midway,below right]{$\vec{v}_1$};
  \draw[->,very thick,vecblue] (0,0)--(-1,1) node[above left]{$\vec{v}_2$};
  % v3 as red
  \draw[->,very thick,vecred] (0,0)--(1,5) node[above right]{$\vec{v}_3$};
  % the construction 2v1 then +3v2
  \draw[->,dashed,vecblue!70] (0,0)--(4,2);
  \draw[->,dashed,vecblue!70] (4,2)--(1,5);
  \node[vecblue!70] at (3.5,1.2) {$2\vec{v}_1$};
  \node[vecblue!70] at (3.0,3.7) {$3\vec{v}_2$};
\end{tikzpicture}
\caption{The redundant vector. $\vec{v}_3$ (red) is reached by walking
$2\vec{v}_1$ then $3\vec{v}_2$ (dashed). It lies inside the reach of the other
two, so the set $\{\vec{v}_1,\vec{v}_2,\vec{v}_3\}$ is dependent.}
\label{ch14:fig-redundant}
\end{figure}

There is a slicker, perfectly symmetric way to say ``one of these is redundant''
that does not single out which one. Move everything to one side of the equation
$\vec{v}_3 = 2\vec{v}_1 + 3\vec{v}_2$:
\[
  2\,\vec{v}_1 + 3\,\vec{v}_2 - \vec{v}_3 = \vec{0}.
\]
We have written the zero vector as a combination of $\vec{v}_1,\vec{v}_2,
\vec{v}_3$ \emph{without using all-zero coefficients}. That is the fingerprint of
dependence, and it is the definition we will lean on for the rest of the book.

\begin{focusbox}
Dependence means the zero vector can be reached in a \emph{nontrivial} way ---
with coefficients that are not all zero. Independence means the \emph{only} way
to reach $\vec{0}$ is the boring way, with every coefficient zero. Everything in
this chapter flows from that one idea.
\end{focusbox}

\section{Linear independence, defined precisely}
\label{ch14:def}

\begin{definition}[Linear independence]
\label{ch14:def-indep}
A list of vectors $\vec{v}_1, \vec{v}_2, \dots, \vec{v}_k$ in a vector space $V$
is \textbf{linearly independent} if the only scalars $c_1, c_2, \dots, c_k$
satisfying
\[
  c_1 \vec{v}_1 + c_2 \vec{v}_2 + \cdots + c_k \vec{v}_k = \vec{0}
\]
are $c_1 = c_2 = \cdots = c_k = 0$. If \emph{any} choice of scalars that are not
all zero produces $\vec{0}$, the list is \textbf{linearly dependent}.
\end{definition}

This definition looks like it is about the zero vector, but it is secretly about
\emph{uniqueness of representation}, which is what we actually care about.
Suppose some vector $\vec{w}$ in the span could be written two different ways,
\[
  \vec{w} = a_1\vec{v}_1 + \cdots + a_k\vec{v}_k
        = b_1\vec{v}_1 + \cdots + b_k\vec{v}_k .
\]
Subtract: $(a_1-b_1)\vec{v}_1 + \cdots + (a_k-b_k)\vec{v}_k = \vec{0}$. If the
$\vec{v}_i$ are independent, every difference $a_i - b_i$ must be zero, so the
two ``different'' representations were the same all along. Independence is
exactly the condition that \emph{coordinates are unambiguous}. Hold that thought;
it is why bases are useful at all.

\begin{trapbox}
``Linearly independent'' is a property of the \emph{whole list together}, not of
individual vectors. No single vector is ``independent.'' And pairwise checks are
not enough: three vectors can be independent two-at-a-time yet dependent as a
trio. In $\R^2$,
\[
  \colvec{1\\0},\ \colvec{0\\1},\ \colvec{1\\1}
\]
has no two parallel vectors, yet $\colvec{1\\0}+\colvec{0\\1}-\colvec{1\\1}=\vec{0}$,
so the three are dependent. Always test the full set at once.
\end{trapbox}

\subsection*{Two vectors: independence is non-parallelism}

For just two vectors the definition collapses to something you can see.
$\{\vec{v}_1, \vec{v}_2\}$ is dependent exactly when one is a scalar multiple of
the other --- when they are parallel (including the degenerate cases where one is
$\vec{0}$). If $c_1\vec{v}_1 + c_2\vec{v}_2 = \vec{0}$ with, say, $c_1 \ne 0$,
then $\vec{v}_1 = -(c_2/c_1)\vec{v}_2$, a scalar multiple. So:

\begin{focusbox}
Two nonzero vectors are independent $\iff$ they are not parallel $\iff$ they span
a genuine plane (not just a line). Independence is the algebra of ``these point
in really different directions.''
\end{focusbox}

\subsection*{The zero vector poisons every set}

If $\vec{v}_1 = \vec{0}$, then $1\cdot\vec{v}_1 + 0\cdot\vec{v}_2 + \cdots +
0\cdot\vec{v}_k = \vec{0}$ uses a nonzero coefficient (the leading $1$), so the
set is automatically dependent. Any list containing $\vec{0}$ is dependent. This
matches the intuition: the zero vector is a road of length zero --- it reaches
nowhere new, so it is always redundant.

\section{The independence test you will actually run}
\label{ch14:test}

Definitions are lovely, but how do you \emph{test} a concrete list? Stack the
vectors as the columns of a matrix and ask elimination to do the work.

The equation $c_1\vec{v}_1 + \cdots + c_k\vec{v}_k = \vec{0}$ is precisely the
matrix equation $\mat{A}\vec{c} = \vec{0}$, where the columns of $\mat{A}$ are the
$\vec{v}_i$ and $\vec{c} = (c_1,\dots,c_k)$. So:

\begin{formalbox}
\textbf{Independence as a nullspace condition.} Let $\mat{A}$ have columns
$\vec{v}_1,\dots,\vec{v}_k$. The columns are linearly independent if and only if
$\nullsp(\mat{A}) = \{\vec{0}\}$, i.e. $\mat{A}\vec{c}=\vec{0}$ has \emph{only}
the trivial solution. Equivalently, after elimination, \emph{every column is a
pivot column} --- there are no free variables.
\end{formalbox}

A free variable is exactly a free dial you can turn to manufacture a nonzero
solution of $\mat{A}\vec{c}=\vec{0}$, i.e. a nontrivial dependence. No free
variables means no slack means independence. Let us run it.

\begin{example}[Testing three vectors in $\R^3$]
\label{ch14:ex-test3}
Are these independent?
\[
  \vec{v}_1=\colvec{1\\2\\3},\quad
  \vec{v}_2=\colvec{1\\1\\1},\quad
  \vec{v}_3=\colvec{3\\5\\7}.
\]
Stack them as columns and eliminate:
\[
  \mat{A}=\begin{bmatrix} 1 & 1 & 3\\ 2 & 1 & 5\\ 3 & 1 & 7\end{bmatrix}
  \xrightarrow[R_3-3R_1]{R_2-2R_1}
  \begin{bmatrix} 1 & 1 & 3\\ 0 & -1 & -1\\ 0 & -2 & -2\end{bmatrix}
  \xrightarrow{R_3-2R_2}
  \begin{bmatrix} 1 & 1 & 3\\ 0 & -1 & -1\\ 0 & 0 & 0\end{bmatrix}.
\]
The third column has no pivot --- column $3$ is a free column. So there is a
nontrivial solution and the vectors are \textbf{dependent}. Read the dependence
off: back-substitute with the free variable $c_3 = 1$. Row $2$ gives $-c_2 - c_3
= 0 \Rightarrow c_2 = -1$. Row $1$ gives $c_1 + c_2 + 3c_3 = 0 \Rightarrow c_1 =
-(-1) - 3 = -2$. Check:
\[
  -2\vec{v}_1 - \vec{v}_2 + \vec{v}_3
  = -2\colvec{1\\2\\3} - \colvec{1\\1\\1} + \colvec{3\\5\\7}
  = \colvec{0\\0\\0}. \checkmark
\]
Indeed $\vec{v}_3 = 2\vec{v}_1 + \vec{v}_2$.
\end{example}

\begin{example}[A genuinely independent trio]
\label{ch14:ex-indep3}
Change one entry:
\[
  \mat{A}=\begin{bmatrix} 1 & 1 & 3\\ 2 & 1 & 5\\ 3 & 1 & 8\end{bmatrix}
  \xrightarrow[R_3-3R_1]{R_2-2R_1}
  \begin{bmatrix} 1 & 1 & 3\\ 0 & -1 & -1\\ 0 & -2 & -1\end{bmatrix}
  \xrightarrow{R_3-2R_2}
  \begin{bmatrix} 1 & 1 & 3\\ 0 & -1 & -1\\ 0 & 0 & 1\end{bmatrix}.
\]
Now there are three pivots, one in every column. No free variables, so
$\mat{A}\vec{c}=\vec{0}$ forces $\vec{c}=\vec{0}$: the columns are
\textbf{independent}. Three independent vectors in $\R^3$ --- as we will see,
that is a basis, and this $\mat{A}$ is invertible.
\end{example}

\begin{labbox}
Test $\colvec{1\\3},\ \colvec{2\\6},\ \colvec{0\\1}$ by hand. (i) Are the first
two independent? (ii) Is the whole set? (iii) You have three vectors in $\R^2$ ---
predict the answer \emph{before} eliminating, then confirm. What general rule
about ``more vectors than dimensions'' are you discovering?
\end{labbox}

\subsection*{More vectors than components: dependence is forced}

That lab points at a theorem you will use constantly. If you have more vectors
than the space has components, redundancy is unavoidable.

\begin{theorem}[Too many vectors must be dependent]
\label{ch14:thm-toomany}
Any list of $k$ vectors in $\R^n$ with $k > n$ is linearly dependent.
\end{theorem}

\begin{proof}
Form the $n \times k$ matrix $\mat{A}$ whose columns are the vectors. The
equation $\mat{A}\vec{c} = \vec{0}$ is a homogeneous system with $n$ equations
and $k$ unknowns. Elimination produces at most $n$ pivots (one per row), so at
most $n$ of the $k$ columns are pivot columns. Since $k > n$, at least one column
is free, giving a free variable and hence a nonzero solution $\vec{c}$. That is a
nontrivial dependence.
\end{proof}

\begin{intuitionbox}
You cannot fit five genuinely independent directions inside three-dimensional
space; by the time you have laid down three independent arrows, the fourth must
live in the room they already span. ``Independent'' caps out at the dimension of
the room. This single fact is the engine behind ``dimension is well-defined,''
which we prove below.
\end{intuitionbox}

\section{Spanning sets: reaching everything}
\label{ch14:span}

Independence is about having no \emph{excess}. Spanning is the opposite concern:
having no \emph{shortage}. We met span in Chapter~3; let us restate it with our
current vocabulary.

\begin{definition}[Spanning set]
\label{ch14:def-span}
Vectors $\vec{v}_1,\dots,\vec{v}_k$ \textbf{span} a vector space $V$ if every
$\vec{w}\in V$ can be written as some linear combination
$\vec{w} = c_1\vec{v}_1 + \cdots + c_k\vec{v}_k$. We write $V = \spn\{\vec{v}_1,
\dots,\vec{v}_k\}$ and call $\{\vec{v}_1,\dots,\vec{v}_k\}$ a spanning set for $V$.
\end{definition}

Spanning says: \emph{no destination in $V$ is out of reach}. In matrix language,
$\vec{v}_1,\dots,\vec{v}_k$ span $\R^n$ if and only if $\mat{A}\vec{x}=\vec{b}$
has a solution for \emph{every} $\vec{b}\in\R^n$ --- equivalently, $\col(\mat{A})
= \R^n$, which (Chapter~12) happens exactly when $\mat{A}$ has a pivot in every
\emph{row}.

\begin{focusbox}
Two complementary questions, two complementary pivot patterns:
\begin{itemize}
  \item \textbf{Independence} (no excess) $\iff$ a pivot in every \emph{column}
        (no free variables).
  \item \textbf{Spanning} $\R^n$ (no shortage) $\iff$ a pivot in every
        \emph{row} (no unsolvable $\vec{b}$).
\end{itemize}
A basis demands both at once: a pivot in every row \emph{and} every column ---
which only a \emph{square}, full-pivot matrix can deliver.
\end{focusbox}

\begin{example}[Spanning but not independent]
\label{ch14:ex-span-not-indep}
In $\R^2$ the three vectors $\colvec{1\\0},\ \colvec{0\\1},\ \colvec{1\\1}$ span
$\R^2$ --- the first two alone already reach everything. But we saw they are
dependent. So you can span with too many vectors; there is fat to trim.
Conversely $\colvec{1\\0}$ alone is independent but spans only the $x$-axis, not
all of $\R^2$ --- independent, but too few. A basis is what you get when you trim
the fat \emph{and} fill the gaps.
\end{example}

The two pivot conditions have a vivid geometric shape if you picture the matrix
$\mat{A}$ as a rectangle of pivots. ``Independent'' is a statement about
\emph{columns}; ``spanning $\R^n$'' is a statement about \emph{rows}; and a basis
is the lucky square where both hold at once.

\begin{figure}[H]
\centering
\begin{tikzpicture}[scale=0.62, font=\small]
  % --- Tall: more rows than columns: can be independent, cannot span ---
  \begin{scope}[xshift=0cm]
    \fill[vecblue!12] (0,0) rectangle (2,4);
    \draw[vecblue!70,thick] (0,0) rectangle (2,4);
    % pivots down the diagonal of the 2 columns
    \fill[vecblue] (0.5,3.5) circle (5pt);
    \fill[vecblue] (1.5,2.5) circle (5pt);
    \node at (1,-0.7) {tall: $m>n$};
    \node[align=center, text width=3.4cm] at (1,-1.9)
      {pivot in every \emph{column}\\ \textbf{independent}, cannot span};
  \end{scope}
  % --- Wide: more columns than rows: can span, cannot be independent ---
  \begin{scope}[xshift=6.4cm]
    \fill[targetgreen!12] (0,0) rectangle (4,2);
    \draw[targetgreen!70,thick] (0,0) rectangle (4,2);
    \fill[targetgreen] (0.5,1.5) circle (5pt);
    \fill[targetgreen] (1.5,0.5) circle (5pt);
    \node at (2,-0.7) {wide: $n>m$};
    \node[align=center, text width=3.8cm] at (2,-1.9)
      {pivot in every \emph{row}\\ \textbf{spans}, cannot be independent};
  \end{scope}
  % --- Square full-rank: basis ---
  \begin{scope}[xshift=13.2cm]
    \fill[structurepurple!14] (0,0) rectangle (3,3);
    \draw[structurepurple!75,thick] (0,0) rectangle (3,3);
    \fill[structurepurple] (0.5,2.5) circle (5pt);
    \fill[structurepurple] (1.5,1.5) circle (5pt);
    \fill[structurepurple] (2.5,0.5) circle (5pt);
    \node at (1.5,-0.7) {square: $m=n$};
    \node[align=center, text width=3.6cm] at (1.5,-1.9)
      {pivot in every row \emph{and} column\\ \textbf{a basis}};
  \end{scope}
\end{tikzpicture}
\caption{The shape of the pivots. Dots mark pivots. A \emph{tall} matrix can have
a pivot in every column (its columns are independent) but never reaches every
$\vec{b}$. A \emph{wide} matrix can have a pivot in every row (its columns span)
but always carries a free column, so its columns are dependent. Only a
\emph{square}, full-pivot matrix achieves a pivot in every row \emph{and} column
at once --- its columns are a basis.}
\label{ch14:fig-pivotshape}
\end{figure}

\section{Basis: independent and spanning}
\label{ch14:basis}

\begin{definition}[Basis]
\label{ch14:def-basis}
A \textbf{basis} of a vector space $V$ is a list of vectors that is
\begin{enumerate}[label=(\roman*)]
  \item linearly independent, and
  \item spans $V$.
\end{enumerate}
\end{definition}

Independent so there is nothing wasted; spanning so there is nothing missing. A
basis is the smallest set that still reaches everything and the largest
independent set that still fits --- those two extremal descriptions turn out to
coincide, which is one of the quiet miracles of the subject.

\begin{figure}[H]
\centering
\begin{tikzpicture}[scale=1.05,>=Stealth]
  % Container box
  \draw[fill=gray!2, draw=gray!25, rounded corners=3pt] (-4.6,-2.4) rectangle (4.6,2.6);
  \node[draw=spacegray!60, fill=gray!10, rounded corners=2pt, font=\footnotesize\scshape, text=black!85]
    at (0,2.2) {The Goldilocks Property of a Basis};

  % Venn circles
  \fill[sourceblue!14] (-1.1,0) circle (1.75);
  \fill[targetgreen!14] (1.1,0) circle (1.75);
  \begin{scope}
    \clip (-1.1,0) circle (1.75);
    \fill[structurepurple!22] (1.1,0) circle (1.75);
  \end{scope}

  \draw[sourceblue!80,thick] (-1.1,0) circle (1.75);
  \draw[targetgreen!80,thick] (1.1,0) circle (1.75);

  % Labels left
  \node[sourceblue!85!black, font=\small\bfseries, align=center] at (-2.3,0.7) {Independent};
  \node[sourceblue!75!black, font=\scriptsize, align=center] at (-2.3,-0.1) {No redundant\\directions};
  \node[sourceblue!60!black, font=\tiny, align=center] at (-2.3,-0.9) {$\le n$ vectors\\(too few can't span)};

  % Labels right
  \node[targetgreen!85!black, font=\small\bfseries, align=center] at (2.3,0.7) {Spanning};
  \node[targetgreen!75!black, font=\scriptsize, align=center] at (2.3,-0.1) {Reaches all\\of the space};
  \node[targetgreen!60!black, font=\tiny, align=center] at (2.3,-0.9) {$\ge n$ vectors\\(too many has waste)};

  % Intersection (Basis)
  \node[structurepurple!95!black, font=\small\bfseries, align=center] at (0,0.4) {Basis};
  \node[structurepurple!85!black, font=\scriptsize, align=center] at (0,-0.2) {Exactly $n$\\vectors};
  \node[draw=structurepurple!70, fill=structurepurple!15, rounded corners=2pt, font=\tiny\bfseries, text=structurepurple!90!black]
    at (0,-0.9) {Just Right};
\end{tikzpicture}
\caption{A basis is the Goldilocks intersection of two complementary conditions.
Independent sets (left) have zero redundancy but may be too small to span;
spanning sets (right) reach everywhere but may contain redundant vectors.
A basis sits exactly at the boundary: minimally spanning and maximally independent.}
\label{ch14:fig-venn}
\end{figure}

\begin{intuitionbox}
A basis is a \emph{coordinate frame} for the space: a minimal set of independent
directions such that every point has \emph{exactly one} address. Independence
guarantees the address is unique (no two combinations give the same vector);
spanning guarantees every point has an address at all. A basis is the bridge
between abstract vectors and the concrete tuples of numbers you compute with.
\end{intuitionbox}

\begin{example}[The standard basis of $\R^n$]
\label{ch14:ex-standard}
The vectors
\[
  \vec{e}_1=\colvec{1\\0\\0},\quad
  \vec{e}_2=\colvec{0\\1\\0},\quad
  \vec{e}_3=\colvec{0\\0\\1}
\]
form the \textbf{standard basis} of $\R^3$. They span (any $\vec{w} = w_1\vec{e}_1
+ w_2\vec{e}_2 + w_3\vec{e}_3$) and they are independent (stacking them gives the
identity matrix, all pivots). The coordinates of $\vec{w}$ in this basis are just
its ordinary entries --- which is why we usually do not even notice we are using a
basis.
\end{example}

\begin{example}[A non-standard basis of $\R^2$]
\label{ch14:ex-nonstandard}
Consider $\vec{b}_1 = \colvec{1\\1}$, $\vec{b}_2 = \colvec{1\\-1}$. Stacked,
\[
  \mat{B}=\begin{bmatrix} 1 & 1\\ 1 & -1\end{bmatrix}
  \xrightarrow{R_2-R_1}
  \begin{bmatrix} 1 & 1\\ 0 & -2\end{bmatrix},
\]
two pivots in a $2\times 2$ matrix: independent and spanning, hence a basis. It
is a perfectly good coordinate frame for the plane --- just rotated and rescaled
from the standard one. We will read off coordinates in this frame in
Section~\ref{ch14:coords}.
\end{example}

\begin{figure}[H]
\centering
\begin{tikzpicture}[>=Stealth, scale=0.8]
  \begin{scope}
    \draw[help lines,gray!25] (-1.5,-1.5) grid (3.5,3.5);
    \draw[->,gray!60] (-1.5,0)--(3.5,0); \draw[->,gray!60] (0,-1.5)--(0,3.5);
    \draw[->,very thick,vecblue] (0,0)--(1,0) node[below right]{$\vec{e}_1$};
    \draw[->,very thick,vecblue] (0,0)--(0,1) node[left]{$\vec{e}_2$};
    \filldraw[vecred] (2,3) circle (2pt) node[above right]{$\vec{w}$};
    \draw[dashed,vecred!70] (2,0)--(2,3)--(0,3);
    \node at (1,-2.3) {standard frame: $\vec{w}=(2,3)$};
  \end{scope}
  \begin{scope}[xshift=8cm]
    \draw[help lines,gray!25] (-1.5,-1.5) grid (3.5,3.5);
    \draw[->,gray!60] (-1.5,0)--(3.5,0); \draw[->,gray!60] (0,-1.5)--(0,3.5);
    % true basis directions: b1=(1,1), b2=(1,-1)
    \draw[->,very thick,structurepurple] (0,0)--(1,1) node[above]{$\vec{b}_1$};
    \draw[->,very thick,structurepurple] (0,0)--(1,-1) node[below right]{$\vec{b}_2$};
    % the same point w=(2,3) reached as (5/2) b1 + (-1/2) b2
    % (5/2) b1 = (2.5,2.5);  then -(1/2) b2 = (-0.5,0.5) lands at (2,3)
    \draw[dashed,structurepurple!70] (0,0)--(2.5,2.5);
    \draw[dashed,structurepurple!70] (2.5,2.5)--(2,3);
    \node[structurepurple!80!black] at (1.6,2.45) {$\tfrac52\vec{b}_1$};
    \node[structurepurple!80!black] at (2.95,2.85) {$-\tfrac12\vec{b}_2$};
    \filldraw[vecred] (2,3) circle (2pt) node[above right]{$\vec{w}$};
    \node[align=center] at (1,-2.3) {$\vec{b}$ frame:\\ $\vec{w}=(\tfrac52,-\tfrac12)$};
  \end{scope}
\end{tikzpicture}
\caption{The same point $\vec{w}$ lives in one fixed space, but its
\emph{coordinates} depend on the basis you choose. Left: in the standard frame
$\vec{w}$ has address $(2,3)$. Right: in the frame $\vec{b}_1=(1,1)$,
$\vec{b}_2=(1,-1)$, the \emph{identical} point is reached by walking
$\tfrac52\vec{b}_1$ then $-\tfrac12\vec{b}_2$ (dashed), so its address is
$(\tfrac52,-\tfrac12)$ --- the numbers computed in Example~\ref{ch14:ex-coords}.
Changing basis is changing the ruler, not the point.}
\label{ch14:fig-frames}
\end{figure}

\begin{strangbox}
Strang hammers home that a basis must do \emph{both} jobs, and he loves the
phrase ``the columns are independent and they span.'' If you are reading this
\emph{before} Strang: the payoff lands in his ``four subspaces'' lecture, where
he hands you a basis for each of the four subspaces straight off the elimination
of $\mat{A}$. If you are reading \emph{after} Strang: notice that his recipe
``pivot columns of $\mat{A}$ form a basis for the column space'' is exactly
Theorem~\ref{ch14:thm-colbasis} below, and the special solutions forming a basis
for the nullspace is Theorem~\ref{ch14:thm-nullbasis}. This chapter is the
``why'' under those recipes.
\end{strangbox}

\section{Dimension is well-defined}
\label{ch14:dim}

Here is the cornerstone. A space has many different bases (we just saw two for
$\R^2$). But \emph{every basis of a given space has the same number of vectors}.
That common number deserves a name.

\begin{definition}[Dimension]
\label{ch14:def-dim}
The \textbf{dimension} of a vector space $V$, written $\dimn V$, is the number of
vectors in any basis of $V$. By the theorem below this number does not depend on
which basis you pick.
\end{definition}

The whole definition would be nonsense if two bases could have different sizes.
The fact that they cannot is the \emph{Steinitz exchange} idea, and it follows
cleanly from Theorem~\ref{ch14:thm-toomany} once we sharpen it slightly.

\begin{lemma}[Independent $\le$ spanning]
\label{ch14:lem-steinitz}
If $\vec{u}_1,\dots,\vec{u}_p$ are linearly independent in $V$, and
$\vec{w}_1,\dots,\vec{w}_q$ span $V$, then $p \le q$. In words: you can never
have \emph{more} independent vectors than you have spanning vectors.
\end{lemma}

\begin{proof}
Because the $\vec{w}_j$ span $V$, each independent vector is a combination of
them: $\vec{u}_i = \sum_{j=1}^{q} a_{ji}\,\vec{w}_j$. Collect the coefficients
into a $q \times p$ matrix $\mat{A}=(a_{ji})$. Now suppose, for contradiction,
that $p > q$. Then $\mat{A}$ has more columns than rows, so by
Theorem~\ref{ch14:thm-toomany} (applied to its columns) the system
$\mat{A}\vec{c}=\vec{0}$ has a nonzero solution $\vec{c}=(c_1,\dots,c_p)$. Form
the same combination of the $\vec{u}_i$:
\[
  \sum_{i=1}^{p} c_i \vec{u}_i
  = \sum_{i=1}^{p} c_i \sum_{j=1}^{q} a_{ji}\vec{w}_j
  = \sum_{j=1}^{q}\Big(\underbrace{\textstyle\sum_{i=1}^{p} a_{ji} c_i}_{=\,0}\Big)\vec{w}_j
  = \vec{0},
\]
where the inner sum is the $j$-th entry of $\mat{A}\vec{c}=\vec{0}$. So a nonzero
$\vec{c}$ gives a nontrivial dependence among the $\vec{u}_i$ ---
contradicting their independence. Hence $p \le q$.
\end{proof}

\begin{theorem}[Dimension is well-defined]
\label{ch14:thm-dim}
Any two bases of a vector space $V$ have the same number of vectors.
\end{theorem}

\begin{proof}
Let $\mathcal B = \{\vec{b}_1,\dots,\vec{b}_p\}$ and $\mathcal C =
\{\vec{c}_1,\dots,\vec{c}_q\}$ both be bases of $V$. Then $\mathcal B$ is
independent and $\mathcal C$ spans, so Lemma~\ref{ch14:lem-steinitz} gives $p \le
q$. By symmetry (now $\mathcal C$ independent, $\mathcal B$ spanning), $q \le p$.
Therefore $p = q$.
\end{proof}

\begin{aestheticbox}
Pause on how strong this is. A subspace is an infinite cloud of vectors, and
there are infinitely many ways to choose a basis for it --- yet every one of
those choices uses \emph{exactly} the same count. Dimension is an invariant: a
property of the \emph{space itself}, blind to your choice of axes. This is the
first time in the book a continuous, infinite object gets summarized by a single
integer, and that integer turns out to control almost everything --- rank,
solvability, the size of the nullspace, the geometry of projections. Whole fields
(topology, algebraic geometry) are built around finding the ``right'' invariant
integer; in linear algebra it is simply dimension, and it is wonderfully
computable.
\end{aestheticbox}

\begin{historybox}
The idea of dimension seems deceptively obvious in daily life: a line is one-dimensional, a sheet of paper is two-dimensional, the room is three-dimensional. But what gives a mathematical subspace an intrinsic dimension, independent of any coordinate system?

In 1844, a visionary German schoolteacher named Hermann Grassmann published \emph{Die Lineare Ausdehnungslehre} (The Theory of Linear Extension). Grassmann introduced the concepts of linear combination, linear independence, and spanning sets, demonstrating that an $n$-dimensional linear manifold could be studied without any ambient coordinate frame. Tragically, Grassmann's writing was so philosophical, dense, and ahead of its time that the mathematical community ignored it for over thirty years. Only in the late nineteenth century did Giuseppe Peano and Felix Klein recognize Grassmann as the true founder of modern vector algebra.

In 1910, Ernst Steinitz published his landmark memoir \emph{Algebraische Theorie der Körper}, formulating the famous \emph{Steinitz exchange lemma} (Lemma~\ref{ch14:lem-steinitz}). Steinitz proved that no matter how two people choose bases for the same vector space, they must always arrive at the exact same number of vectors. Dimension is not an artifact of human choice; it is an unshakeable topological invariant of the space.

While finite-dimensional dimension was settled by Steinitz, the concept ignited an intense controversy when applied to infinite-dimensional spaces such as the space of continuous functions $C[0,1]$. In 1905, Georg Hamel defined a basis using strictly finite linear combinations (a \emph{Hamel basis}). Proving that every vector space has a Hamel basis requires Zorn's Lemma---and thus the controversial Axiom of Choice. For spaces like $C[0,1]$, a Hamel basis is uncountably infinite, and nobody can ever write down an explicit example. To escape this non-constructive trap, functional analysts led by Juliusz Schauder in 1927 introduced \emph{Schauder bases}, which allow infinite convergent series $\sum_{k=1}^\infty c_k \vec{e}_k$ in normed spaces (like Fourier series $\sin(nx)$ and $\cos(nx)$). The debate between pure algebraists insisting on finite algebraic sums and analysts demanding topological limits reshaped twentieth-century mathematics, demarcating the boundary between linear algebra and functional analysis.
\end{historybox}

\subsection*{Consequences you will reuse constantly}

\begin{corollary}[Counting shortcuts in $\R^n$]
\label{ch14:cor-count}
In an $n$-dimensional space:
\begin{enumerate}[label=(\alph*)]
  \item Any independent list has at most $n$ vectors; more than $n$ is forced
        dependent.
  \item Any spanning list has at least $n$ vectors; fewer than $n$ cannot span.
  \item Any \emph{independent list of exactly $n$ vectors is automatically a
        basis} (it must already span).
  \item Any \emph{spanning list of exactly $n$ vectors is automatically a basis}
        (it must already be independent).
\end{enumerate}
\end{corollary}

Parts (c) and (d) are the great labor-savers: in $\R^n$, once you have the right
\emph{count}, you only need to check \emph{one} of the two conditions. If you have
$n$ vectors and they are independent, you are done --- spanning comes free, and
vice versa.

\begin{example}[Using the count]
\label{ch14:ex-count}
Are $\colvec{2\\1\\0},\colvec{0\\3\\1},\colvec{1\\0\\2}$ a basis of $\R^3$? There
are exactly three of them in a $3$-dimensional space, so by
Corollary~\ref{ch14:cor-count}(c) it is enough to check independence. Eliminate:
\[
  \begin{bmatrix} 2 & 0 & 1\\ 1 & 3 & 0\\ 0 & 1 & 2\end{bmatrix}
  \xrightarrow{R_2-\frac12 R_1}
  \begin{bmatrix} 2 & 0 & 1\\ 0 & 3 & -\frac12\\ 0 & 1 & 2\end{bmatrix}
  \xrightarrow{R_3-\frac13 R_2}
  \begin{bmatrix} 2 & 0 & 1\\ 0 & 3 & -\frac12\\ 0 & 0 & \frac{13}{6}\end{bmatrix}.
\]
Three pivots, so independent, so (no spanning check needed) a basis.
\end{example}

\begin{walkbox}
Take a walk and convince yourself, with no algebra, of each of the four parts of
Corollary~\ref{ch14:cor-count} using only the picture from
Theorem~\ref{ch14:thm-toomany}: ``you cannot cram more independent directions
into a room than the room has dimensions, and you cannot cover the room with
fewer reaching-directions than that.'' Which single inequality from
Lemma~\ref{ch14:lem-steinitz} is doing all the work?
\end{walkbox}

\section{Coordinates in a basis}
\label{ch14:coords}

We said independence makes representations unique. Let us cash that in. Once you
fix a basis, every vector gets a unique list of numbers --- its
\textbf{coordinates} in that basis.

\begin{definition}[Coordinates]
\label{ch14:def-coords}
Let $\mathcal B = \{\vec{b}_1,\dots,\vec{b}_n\}$ be a basis of $V$. For $\vec{w}
\in V$ there are unique scalars $c_1,\dots,c_n$ with $\vec{w} = c_1\vec{b}_1 +
\cdots + c_n\vec{b}_n$. The vector
\[
  [\vec{w}]_{\mathcal B} = \colvec{c_1\\\vdots\\c_n}
\]
is the \textbf{coordinate vector} of $\vec{w}$ with respect to $\mathcal B$.
\end{definition}

Uniqueness is precisely the independence argument from Section~\ref{ch14:def}: two
addresses for the same point would subtract to a nontrivial dependence, which a
basis forbids. So coordinates are well-defined, and finding them is just solving a
square system.

\begin{example}[Coordinates in the frame $\{\vec{b}_1,\vec{b}_2\}$]
\label{ch14:ex-coords}
Use $\vec{b}_1 = \colvec{1\\1}$, $\vec{b}_2 = \colvec{1\\-1}$ from
Example~\ref{ch14:ex-nonstandard}, and find the coordinates of $\vec{w} =
\colvec{2\\3}$. We need $c_1\vec{b}_1 + c_2\vec{b}_2 = \vec{w}$, i.e.
\[
  \begin{bmatrix} 1 & 1\\ 1 & -1\end{bmatrix}\colvec{c_1\\c_2} = \colvec{2\\3}.
\]
Adding the two equations: $2c_1 = 5 \Rightarrow c_1 = \tfrac52$. Subtracting:
$2c_2 = -1 \Rightarrow c_2 = -\tfrac12$. So
\[
  [\vec{w}]_{\mathcal B} = \colvec{\tfrac52\\[2pt]-\tfrac12},
  \qquad\text{meaning}\qquad
  \tfrac52\colvec{1\\1} - \tfrac12\colvec{1\\-1} = \colvec{2\\3}.\ \checkmark
\]
The point $\vec{w}$ never moved. In the standard frame its address is $(2,3)$; in
the $\mathcal B$ frame its address is $(\tfrac52,-\tfrac12)$. Same point, different
ruler.
\end{example}

\begin{focusbox}
To find coordinates, put the basis vectors in the columns of $\mat{B}$ and solve
$\mat{B}\,[\vec{w}]_{\mathcal B} = \vec{w}$. Because $\mathcal B$ is a basis,
$\mat{B}$ is invertible, so $[\vec{w}]_{\mathcal B} = \mat{B}\inv\vec{w}$ and the
answer is unique. ``Change of basis'' is multiplication by $\mat{B}\inv$ ---
the full story waits in the change-of-basis chapter, but the seed is right here.
\end{focusbox}

\begin{connectionbox}
This is the idea that powers \emph{Fourier series} and \emph{wavelets}. There, the
``vectors'' are functions and the basis is a family of sines and cosines (or
little wavelet bumps). The Fourier coefficients of a signal are nothing but its
\emph{coordinates} in that basis --- exactly Definition~\ref{ch14:def-coords},
carried into an infinite-dimensional space. Independence becomes orthogonality,
solving $\mat{B}\vec{c}=\vec{w}$ becomes an integral, and ``every function has a
unique coordinate list'' becomes the completeness of the Fourier basis. The
discrete cosine transform inside every JPEG is this chapter, applied to
$8\times 8$ blocks of pixels.
\end{connectionbox}

\section{A basis for the column space}
\label{ch14:colbasis}

Now we put all of this to work on the two subspaces that matter most. Start with
the column space $\col(\mat{A})$, the span of the columns of $\mat{A}$ ---
everything $\mat{A}$ can output.

The columns span $\col(\mat{A})$ by definition, but they may not be independent
(some columns can be redundant). To get a \emph{basis} we trim the redundant ones.
Elimination tells us exactly which to keep.

\begin{theorem}[Pivot columns are a basis for $\col(\mat{A})$]
\label{ch14:thm-colbasis}
The pivot columns of $\mat{A}$ --- the columns of the \emph{original} matrix in
the positions where its row-reduced form has pivots --- form a basis for
$\col(\mat{A})$. Hence $\dimn \col(\mat{A})$ equals the number of pivots, which is
the rank of $\mat{A}$.
\end{theorem}

\begin{proof}[Why it works]
Row operations do not change which combinations of columns give zero: if $\mat{R}$
is the reduced form of $\mat{A}$, then $\mat{A}\vec{c}=\vec{0}$ and
$\mat{R}\vec{c}=\vec{0}$ have the same solutions $\vec{c}$. In $\mat{R}$ the
pivot columns are clearly independent (each carries a pivot the others lack) and
each free column is visibly a combination of earlier pivot columns. The same
dependence relations $\vec{c}$ hold among the columns of $\mat{A}$ --- so the
pivot columns of $\mat{A}$ are independent, and the free columns of $\mat{A}$ are
combinations of them. Independent and spanning the column space: a basis. (One
caution lives in the trap box below.)
\end{proof}

\begin{example}[Reading a column-space basis]
\label{ch14:ex-colbasis}
Let
\[
  \mat{A}=\begin{bmatrix} 1 & 2 & 1 & 3\\ 2 & 4 & 3 & 7\\ 1 & 2 & 2 & 4\end{bmatrix}.
\]
Eliminate. Subtract $2R_1$ from $R_2$ and $R_1$ from $R_3$:
\[
  \begin{bmatrix} 1 & 2 & 1 & 3\\ 0 & 0 & 1 & 1\\ 0 & 0 & 1 & 1\end{bmatrix}
  \xrightarrow{R_3-R_2}
  \begin{bmatrix} 1 & 2 & 1 & 3\\ 0 & 0 & 1 & 1\\ 0 & 0 & 0 & 0\end{bmatrix}.
\]
Pivots sit in columns $1$ and $3$. So a basis for $\col(\mat{A})$ is
\[
  \left\{\ \colvec{1\\2\\1},\ \colvec{1\\3\\2}\ \right\}
  \qquad(\text{columns } 1 \text{ and } 3 \text{ of the \emph{original} } \mat{A}),
\]
and $\dimn\col(\mat{A}) = \rank(\mat{A}) = 2$. The column space is a plane in
$\R^3$. Columns $2$ and $4$ were redundant: indeed column $2 = 2\times$ column $1$,
and you can check column $4 = 2(\text{col }1) + (\text{col }3)$.
\end{example}

\begin{trapbox}
Take the pivot columns from the \textbf{original} $\mat{A}$, \emph{never} the
columns of the row-reduced $\mat{R}$. Row operations preserve the column space's
\emph{dimension} and the \emph{dependence pattern}, but they scramble the actual
column vectors --- $\col(\mat{R})$ is usually a different subspace from
$\col(\mat{A})$. In Example~\ref{ch14:ex-colbasis}, $\mat{R}$'s columns all have a
$0$ in the last slot, so $\col(\mat{R})$ is the $xy$-plane; but $\col(\mat{A})$ is
a \emph{tilted} plane. Same dimension, different plane. Use $\mat{A}$'s columns.
\end{trapbox}

Here is what that ``same dimension, different plane'' caution looks like.
Both planes are $2$-dimensional, but they sit in $\R^3$ at different tilts: the
column space of $\mat{A}$ is spanned by the two genuine pivot columns
$\colvec{1\\2\\1}$ and $\colvec{1\\3\\2}$, while the row-reduced columns flatten
onto the floor.

\begin{figure}[H]
\centering
\begin{tikzpicture}[scale=2.0, line join=round, >=Stealth]
  % axes
  \draw[->,gray!60] (0,0,0)--(1.7,0,0) node[right]{$x$};
  \draw[->,gray!60] (0,0,0)--(0,1.6,0) node[above]{$z$};
  \draw[->,gray!60] (0,0,0)--(0,0,1.7) node[below left]{$y$};
  % the flat xy-plane = col(R): the floor patch
  \fill[gray!18, opacity=0.7] (0,0,0)--(1.3,0,0)--(1.3,0,1.3)--(0,0,1.3)--cycle;
  \node[gray!60!black] at (1.15,0,1.15) {$\col(\mat{R})$};
  % the tilted col(A) plane: spanned by p1=(1,2,1) and p2=(1,3,2)
  % use directions a=(1,1,2)/|.| and b=(1,2,1)/|.| scaled small for a parallelogram
  % Build patch from origin with two spanning vectors u=(0.6,0.55,0.3), v=(0.3,0.7,0.55)
  \coordinate (O) at (0,0,0);
  \fill[structurepurple!22, opacity=0.75]
    (0,0,0) -- (0.6,0.55,0.3) -- (0.9,1.25,0.85) -- (0.3,0.7,0.55) -- cycle;
  \draw[structurepurple!75,thick]
    (0,0,0) -- (0.6,0.55,0.3) -- (0.9,1.25,0.85) -- (0.3,0.7,0.55) -- cycle;
  % the two pivot-column basis vectors drawn inside the tilted plane
  \draw[->,very thick,vecblue] (0,0,0)--(0.6,0.55,0.3)
        node[below right]{col 1};
  \draw[->,very thick,vecred] (0,0,0)--(0.3,0.7,0.55)
        node[above left]{col 3};
  \node[structurepurple!75!black] at (0.95,1.15,0.6) {$\col(\mat{A})$};
\end{tikzpicture}
\caption{Two planes of the \emph{same} dimension in $\R^3$. The row-reduced
matrix $\mat{R}$ has all its columns lying flat in the $xy$-plane (gray floor),
so $\col(\mat{R})$ is the floor. The original $\mat{A}$ spans a \emph{tilted}
plane (purple) carried by its two true pivot columns (blue and red). Equal
dimension, different subspace --- which is why a column-space basis must come
from $\mat{A}$, not from $\mat{R}$.}
\label{ch14:fig-tilted}
\end{figure}

\section{A basis for the nullspace}
\label{ch14:nullbasis}

The nullspace $\nullsp(\mat{A})$ is the set of all $\vec{x}$ with $\mat{A}\vec{x}
= \vec{0}$ --- every combination of columns that cancels to zero, the ``recipes
for nothing.'' Its basis comes from the \textbf{special solutions}: set one free
variable to $1$ and the rest to $0$, then solve for the pivot variables.

\begin{theorem}[Special solutions are a basis for $\nullsp(\mat{A})$]
\label{ch14:thm-nullbasis}
If $\mat{A}$ has $r$ pivots and $f$ free columns (so the number of columns is $n =
r + f$), then the $f$ special solutions form a basis for $\nullsp(\mat{A})$. Hence
\[
  \dimn\nullsp(\mat{A}) = f = n - r = (\#\text{columns}) - \rank(\mat{A}).
\]
\end{theorem}

\begin{proof}[Why it works]
\emph{They span.} Every solution of $\mat{A}\vec{x}=\vec{0}$ is determined by its
free-variable values, and the general solution is exactly the combination of the
special solutions weighted by those free values --- that is how we built them.
\emph{They are independent.} Look at the free-variable slots: special solution $i$
has a $1$ in slot $i$ and $0$ in the other free slots. If a combination
$\sum_i c_i\vec{s}_i = \vec{0}$, then reading off free slot $i$ gives $c_i = 0$.
So all $c_i$ vanish: independent. Independent and spanning the nullspace: a basis.
\end{proof}

The pattern of $1$'s and $0$'s in the free slots is the same trick that made the
standard basis independent --- it is independence handed to you for free by
construction.

\begin{example}[Reading a nullspace basis]
\label{ch14:ex-nullbasis}
Reuse $\mat{A}$ from Example~\ref{ch14:ex-colbasis}, whose reduced form (after one
more cleanup, $R_1 - R_2$, to clear the pivot column) is
\[
  \mat{R}=\begin{bmatrix} 1 & 2 & 0 & 2\\ 0 & 0 & 1 & 1\\ 0 & 0 & 0 & 0\end{bmatrix}.
\]
Pivot variables: $x_1, x_3$. Free variables: $x_2, x_4$. The equations are
\[
  x_1 + 2x_2 + 2x_4 = 0,\qquad x_3 + x_4 = 0.
\]
\textbf{Special solution 1} ($x_2=1,\ x_4=0$): $x_3 = 0$, $x_1 = -2$, giving
$\vec{s}_1 = (-2,1,0,0)$.
\textbf{Special solution 2} ($x_2=0,\ x_4=1$): $x_3 = -1$, $x_1 = -2$, giving
$\vec{s}_2 = (-2,0,-1,1)$. So
\[
  \nullsp(\mat{A}) = \spn\!\left\{\ \colvec{-2\\1\\0\\0},\ \colvec{-2\\0\\-1\\1}\ \right\},
  \qquad \dimn\nullsp(\mat{A}) = 2.
\]
Sanity check the bookkeeping: $n = 4$ columns, $r = 2$ pivots, so $f = n - r = 2$
free variables --- exactly the two basis vectors we found. And $2 + 2 = 4$: the
dimension of the column space plus the dimension of the nullspace equals the
number of columns. That balance is the \emph{Rank--Nullity Theorem}, the headline
act of the next chapter; you are seeing its seed here.
\end{example}

\begin{applicationbox}
\textbf{Why data scientists count dimensions.} A data matrix $\mat{X}$ stores one
example per row and one feature per column. If two features are linearly dependent
--- say \texttt{height\_cm} and \texttt{height\_in}, or
\texttt{revenue}, \texttt{cost}, and \texttt{profit} $=$ \texttt{revenue} $-$
\texttt{cost} --- then those columns are redundant and $\rank(\mat{X})$ is
strictly less than the number of features. In a linear regression this redundancy
makes the normal equations \emph{singular}: there is a whole nullspace of weight
vectors that fit the data equally well, so the ``best coefficients'' are not
unique and the fit becomes numerically unstable. The cure is exactly this
chapter: find a \emph{basis} for the column space of $\mat{X}$ (a maximal
independent set of features), keep those, and drop the rest. ``Effective number
of independent features'' is literally $\rank(\mat{X})$, and the redundant
directions you delete are spanned by the special solutions in $\nullsp(\mat{X})$.
The same dimension count reappears as the number of nonzero singular values in
the SVD, and as the number of principal components worth keeping in PCA --- whole
later chapters are this idea wearing fancier clothes.
\end{applicationbox}

\section{Bases in other vector spaces}
\label{ch14:other}

Nothing in Definition~\ref{ch14:def-basis} mentioned $\R^n$. Independence,
spanning, basis, and dimension live in \emph{any} vector space, and that
generality is where the concepts earn their keep.

\begin{example}[Polynomials of degree $\le 2$]
\label{ch14:ex-poly}
Let $V = \mathcal P_2$, the polynomials $a + bx + cx^2$. A natural basis is
$\{1,\, x,\, x^2\}$: any polynomial is a unique combination of them (spanning $+$
independence), so $\dimn \mathcal P_2 = 3$. The coordinate vector of $3 - 5x +
2x^2$ in this basis is $(3,-5,2)$. Different basis, different coordinates: in the
basis $\{1,\, x-1,\, (x-1)^2\}$ (a Taylor-at-$1$ frame), the \emph{same}
polynomial gets a different coordinate triple, but it is still a $3$-dimensional
space. Once you have a basis, $\mathcal P_2$ behaves \emph{exactly} like $\R^3$ ---
that quiet equivalence (an isomorphism) is the deepest payoff of having a basis.
\end{example}

\begin{example}[A two-dimensional space of matrices]
\label{ch14:ex-matrices}
The symmetric $2\times 2$ matrices with trace zero, $\begin{bmatrix} a & b\\ b &
-a\end{bmatrix}$, form a vector space. A basis is
\[
  \begin{bmatrix} 1 & 0\\ 0 & -1\end{bmatrix},\qquad
  \begin{bmatrix} 0 & 1\\ 1 & 0\end{bmatrix},
\]
since every such matrix is $a$ times the first plus $b$ times the second, uniquely.
Dimension $2$. (Physics readers: these are two of the Pauli matrices, and this
$2$-dimensional space is where qubit observables live.)
\end{example}

\begin{focusbox}
The machinery is universal. Whenever you can add and scale --- arrows, polynomials,
matrices, functions, signals --- you can ask ``what is a basis?'' and ``what is the
dimension?'' The \emph{answers} differ, but the \emph{questions}, and the
elimination-style tools for answering them, are always the ones in this chapter.
\end{focusbox}

\section{Putting it together}
\label{ch14:together}

\begin{recoverybox}
A quick self-check before the exercises. Make sure you can say, in one breath:
\begin{itemize}
  \item \textbf{Independent} $=$ no redundant vector $=$ only the trivial
        combination gives $\vec{0}$ $=$ a pivot in every column $=$
        $\nullsp = \{\vec{0}\}$.
  \item \textbf{Spans $\R^n$} $=$ reaches every $\vec{b}$ $=$ a pivot in every row.
  \item \textbf{Basis} $=$ independent $+$ spanning $=$ a minimal reaching set $=$
        a maximal fitting set $=$ unique coordinates for every vector.
  \item \textbf{Dimension} $=$ size of any basis (well-defined!) $=$ number of
        pivots, for a column space.
  \item \textbf{Column-space basis} $=$ pivot columns of the \emph{original}
        $\mat{A}$. \textbf{Nullspace basis} $=$ the special solutions.
\end{itemize}
\end{recoverybox}

Everything fits one picture. Elimination on $\mat{A}$ splits its columns into
pivot columns and free columns. The pivot columns, taken from the original
matrix, are a basis for the column space; their count is the rank. The free
columns each donate one special solution, and those are a basis for the
nullspace; their count is $n - r$. Two subspaces, two bases, two dimensions ---
all falling out of a single row reduction. In the next chapter those two
dimensions, plus the two from $\mat{A}\T$, lock together into the Four
Fundamental Subspaces and the Rank--Nullity Theorem.

\begin{connectionbox}
We can now determine the exact intrinsic size of any subspace ($\dim V$) and guarantee that every vector has a unique coordinate representation relative to any chosen basis. The pivot columns of $\mat{A}$ give an explicit basis for $\col(\mat{A})$, and the special solutions give an explicit basis for $\nullsp(\mat{A})$.

However, a general basis can be arbitrarily crooked and skewed. If two basis vectors point in nearly the same direction, computing coordinates requires inverting a horribly ill-conditioned matrix. Furthermore, a general basis provides no concept of perpendicularity, distance, or optimal geometric projection. In Chapters 15 and 16, we discover how the dimensions of the four fundamental subspaces lock together into the Fundamental Theorem of Linear Algebra. Then, in Part~IV (Chapters 17--20), we cure the disease of crooked bases by constructing orthonormal bases via the Gram--Schmidt process and QR factorization.
\end{connectionbox}

\section*{Exercises}
\addcontentsline{toc}{section}{Exercises}

\begin{exercise}[Warmup: parallel or not]
\label{ch14:exr-parallel}
Decide whether each pair is independent. (a) $\colvec{2\\-1},\colvec{-6\\3}$.
(b) $\colvec{1\\2},\colvec{2\\1}$. (c) $\colvec{0\\0},\colvec{4\\7}$. For each,
state the reason in one sentence using the word ``parallel'' or ``zero.''
\end{exercise}

\begin{exercise}[Run the test]
\label{ch14:exr-test}
Use elimination to decide whether $\colvec{1\\0\\2},\colvec{2\\1\\1},
\colvec{4\\1\\5}$ are independent. If they are dependent, find an explicit
nontrivial combination equal to $\vec{0}$ and verify it.
\end{exercise}

\begin{exercise}[Count first, then check]
\label{ch14:exr-count}
You are handed four vectors in $\R^3$. Without any computation, what can you
conclude about their independence, and which theorem says so? Now: are five
vectors enough to \emph{span} $\R^3$? Could three? Could two? Justify each with
Corollary~\ref{ch14:cor-count}.
\end{exercise}

\begin{exercise}[Build a basis by trimming]
\label{ch14:exr-trim}
The set $\colvec{1\\1\\0},\colvec{2\\2\\0},\colvec{0\\1\\1},\colvec{1\\0\\-1}$
spans some subspace of $\R^3$. Find a basis for that subspace by selecting pivot
columns, and state its dimension. Is the subspace all of $\R^3$, a plane, or a
line?
\end{exercise}

\begin{exercise}[Coordinates]
\label{ch14:exr-coords}
With the basis $\vec{b}_1=\colvec{2\\1},\vec{b}_2=\colvec{1\\3}$ of $\R^2$, find
$[\vec{w}]_{\mathcal B}$ for $\vec{w}=\colvec{5\\5}$. Then find the vector whose
$\mathcal B$-coordinates are $(1,-1)$. Sketch both in the $\mathcal B$ frame.
\end{exercise}

\begin{exercise}[Column space and nullspace together]
\label{ch14:exr-both}
For
\[
  \mat{A}=\begin{bmatrix} 1 & 3 & 1 & 4\\ 2 & 6 & 1 & 7\\ 1 & 3 & 0 & 3\end{bmatrix},
\]
find a basis for $\col(\mat{A})$ and a basis for $\nullsp(\mat{A})$. Verify that
$\dimn\col(\mat{A}) + \dimn\nullsp(\mat{A})$ equals the number of columns.
\end{exercise}

\begin{exercise}[A basis outside $\R^n$]
\label{ch14:exr-poly}
Show that $\{1+x,\ 1-x,\ x^2\}$ is a basis of $\mathcal P_2$ (polynomials of
degree $\le 2$). Then find the coordinates of $p(x) = 4 + 2x + 3x^2$ in this
basis.
\end{exercise}

\begin{exercise}[Extend to a basis]
\label{ch14:exr-extend}
The vectors $\colvec{1\\1\\1}$ and $\colvec{1\\0\\-1}$ are independent in $\R^3$.
Find a third vector that makes them a basis of $\R^3$, and prove your choice
works. (Hint: you only need to add one direction not already in their span; the
count then guarantees a basis.)
\end{exercise}

\begin{exercise}[Proof: a spanning set can be shrunk to a basis]
\label{ch14:exr-shrink}
Suppose $\vec{v}_1,\dots,\vec{v}_k$ span a space $V$ but are dependent. Prove that
some $\vec{v}_i$ is a combination of the others, and that deleting it leaves a set
that still spans $V$. (This is the trimming process; iterating it reaches a
basis.)
\end{exercise}

\begin{exercise}[Proof: independence and uniqueness are the same thing]
\label{ch14:exr-unique}
Prove that $\vec{v}_1,\dots,\vec{v}_k$ are linearly independent if and only if
every vector in $\spn\{\vec{v}_1,\dots,\vec{v}_k\}$ has \emph{exactly one}
expression as a combination $c_1\vec{v}_1+\cdots+c_k\vec{v}_k$. (One direction is
the argument in Section~\ref{ch14:def}; supply both.)
\end{exercise}

\begin{exercise}[Challenge: dimension by two routes]
\label{ch14:exr-challenge}
Let $W = \spn\{\vec{u}_1,\vec{u}_2,\vec{u}_3\}$ where $\vec{u}_1=\colvec{1\\2\\0\\1},
\vec{u}_2=\colvec{2\\4\\1\\3},\vec{u}_3=\colvec{0\\0\\1\\1}$ in $\R^4$. (a) Find
$\dimn W$ by selecting pivot columns. (b) Independently, write $W$ as the
nullspace of some matrix --- i.e. find the linear equations cutting out $W$ ---
and recompute $\dimn W$ from that description as $4 - (\text{number of independent
equations})$. Confirm the two answers agree, and explain in a sentence why they
must.
\end{exercise}
